\documentclass[../DoAn.tex]{subfiles}
\begin{document}

Chương này giới thiệu các công nghệ và nền tảng lý thuyết được sử dụng trong hệ thống. Với mỗi công nghệ, chương trình bày nguyên lý hoạt động cơ bản, liệt kê các lựa chọn thay thế, và phân tích lý do lựa chọn dựa trên các yêu cầu đã xác định ở Chương 2.

\section{Phần cứng nhúng}
\label{section:3.1}

\subsection{Vi điều khiển ESP32}
\label{subsection:3.1.1}

ESP32 là vi điều khiển 32-bit do Espressif Systems phát triển, sử dụng lõi Xtensa LX6 dual-core với xung nhịp tối đa 240\,MHz \cite{espressif_esp32}. Chip tích hợp sẵn Wi-Fi chuẩn IEEE 802.11 b/g/n, Bluetooth 4.2 và Bluetooth Low Energy (BLE), cùng 520\,KB SRAM và hỗ trợ bộ nhớ Flash ngoài lên đến 16\,MB. Ngoài ra, ESP32 cung cấp đầy đủ các giao tiếp ngoại vi như UART, SPI, I\textsuperscript{2}C và ADC, cho phép kết nối linh hoạt với nhiều loại cảm biến và module ngoại vi. Trong hệ thống này, giao tiếp UART được sử dụng để trao đổi dữ liệu với module SIM7600.

Đặc điểm nổi bật của ESP32 đối với hệ thống đề xuất là khả năng giao tiếp đa giao thức và xử lý đồng thời nhiều tác vụ. ESP32 sử dụng UART2 để kết nối với module SIM7600, nhận dữ liệu GPS và điều khiển modem thông qua tập lệnh AT command được thư viện TinyGSM đóng gói thành các API mức cao. Bên cạnh đó, ESP32 hỗ trợ ngắt ngoài (External Interrupt) trên hầu hết các chân GPIO, cho phép xử lý sự kiện nhấn nút SOS ngay khi xảy ra mà không cần liên tục kiểm tra trạng thái bằng phương pháp polling, từ đó giảm tải cho bộ xử lý và tăng khả năng đáp ứng của hệ thống.

Các phương án thay thế đã được xem xét gồm: Arduino Uno (tài nguyên xử lý và bộ nhớ hạn chế, không tích hợp kết nối mạng không dây và khó đáp ứng yêu cầu truyền dữ liệu bảo mật TLS), Raspberry Pi Zero W (sử dụng hệ điều hành Linux đầy đủ, dẫn đến mức tiêu thụ điện năng và thời gian khởi động cao hơn, không phù hợp với các thiết bị IoT hoạt động liên tục bằng nguồn pin), và STM32 (có hiệu năng tốt nhưng việc phát triển ứng dụng sử dụng MQTT, TinyGSM và các thư viện Arduino cần nhiều cấu hình hơn, làm tăng thời gian phát triển). ESP32 được lựa chọn nhờ khả năng cân bằng giữa hiệu năng xử lý, tích hợp kết nối không dây, mức tiêu thụ điện năng thấp trong chế độ tiết kiệm năng lượng, cùng hệ sinh thái Arduino và ESP-IDF phong phú, giúp rút ngắn thời gian phát triển và bảo trì hệ thống.

Trong hệ thống đề xuất, ESP32 đóng vai trò là bộ điều khiển trung tâm, chịu trách nhiệm giao tiếp với module SIM7600, xử lý dữ liệu GPS, giám sát trạng thái nút SOS, đóng gói dữ liệu theo định dạng JSON và truyền dữ liệu đến máy chủ thông qua giao thức MQTT.

\subsection{Module SIM7600}
\label{subsection:3.1.2}

SIM7600 là module truyền thông đa băng tần do SIMCom Wireless Solutions sản xuất, tích hợp đồng thời bộ thu tín hiệu GNSS (GPS/GLONASS/BeiDou) và modem di động hỗ trợ mạng 4G LTE Cat-1, đồng thời có khả năng chuyển xuống mạng 3G hoặc 2G khi cần thiết trong cùng một thiết bị \cite{simcom_sim7600}. Module giao tiếp với vi điều khiển thông qua giao diện UART sử dụng tập lệnh AT command tiêu chuẩn, cho phép ESP32 điều khiển cả chức năng định vị và kết nối dữ liệu thông qua một kênh giao tiếp duy nhất.

Về khả năng định vị, SIM7600 tích hợp bộ thu tín hiệu GNSS hỗ trợ nhiều hệ thống vệ tinh như GPS, GLONASS và BeiDou, cho phép xác định vị trí của thiết bị phục vụ chức năng theo dõi trong hệ thống. Module hỗ trợ cả khởi động lạnh (Cold Start) và khởi động nóng (Hot Start), giúp cải thiện thời gian thu nhận tín hiệu vệ tinh khi thông tin định vị đã được lưu trước đó.

Đối với kết nối dữ liệu, SIM7600 hỗ trợ mạng di động 4G LTE và có khả năng tự động chuyển xuống mạng 3G hoặc 2G khi cần thiết. Nhờ đó, thiết bị có thể duy trì kết nối mạng trong nhiều điều kiện phủ sóng khác nhau, đáp ứng yêu cầu truyền dữ liệu định vị và các bản tin MQTT từ thiết bị đến máy chủ một cách ổn định.

Firmware giao tiếp với SIM7600 thông qua thư viện TinyGSM. Thư viện này đóng gói phần lớn các tập lệnh AT command thường dùng thành các API mức cao, chẳng hạn như \texttt{modem.getGPS()} để lấy thông tin vị trí hoặc \texttt{modem.gprsConnect()} để thiết lập kết nối dữ liệu. Nhờ đó, mã nguồn trở nên ngắn gọn, dễ bảo trì và giảm sự phụ thuộc vào cú pháp AT command của từng dòng module.

Các phương án thay thế đã được xem xét gồm: kết hợp module GPS u-blox NEO-6M với modem SIM800L (sử dụng hai module riêng biệt làm tăng số lượng linh kiện, diện tích bo mạch và độ phức tạp trong thiết kế phần cứng), module SIM808 (tích hợp GPS nhưng chỉ hỗ trợ mạng 2G với băng thông thấp), và module SIM7000G (hỗ trợ LTE-M/NB-IoT với mức tiêu thụ điện năng thấp hơn nhưng tốc độ truyền dữ liệu hạn chế và hạ tầng NB-IoT/LTE-M tại Việt Nam hiện chưa được triển khai rộng bằng mạng 4G LTE). Vì vậy, SIM7600 được lựa chọn do tích hợp cả chức năng định vị GPS và kết nối 4G LTE trong một module duy nhất, giúp đơn giản hóa thiết kế phần cứng, giảm độ phức tạp của firmware và đảm bảo khả năng kết nối ổn định trên hạ tầng mạng di động hiện nay.

Trong hệ thống đề xuất, SIM7600 đảm nhiệm đồng thời hai chức năng chính: xác định vị trí GPS của thiết bị và cung cấp kết nối mạng di động để truyền dữ liệu đến máy chủ. Nhờ đó, thiết bị có thể hoạt động độc lập ở mọi nơi có vùng phủ sóng di động mà không phụ thuộc vào mạng Wi-Fi.

\section{Giao thức truyền thông}
\label{section:3.2}

\subsection{Giao thức MQTT}
\label{subsection:3.2.1}

MQTT (Message Queuing Telemetry Transport) là giao thức truyền thông theo mô hình publish/subscribe được thiết kế dành cho các thiết bị có tài nguyên hạn chế và hoạt động trên các mạng có băng thông thấp hoặc không ổn định. Giao thức được chuẩn hóa bởi tổ chức OASIS và hoạt động trên tầng ứng dụng dựa trên giao thức TCP/IP \cite{mqtt_standard}. Với phần tiêu đề có kích thước nhỏ, MQTT giúp giảm lượng dữ liệu truyền trên mạng, phù hợp với các ứng dụng IoT sử dụng kết nối di động.

Kiến trúc MQTT gồm ba thành phần chính là publisher, broker và subscriber. Trong hệ thống đề xuất, thiết bị ESP32 đóng vai trò là publisher, thực hiện gửi dữ liệu vị trí GPS và tín hiệu cảnh báo SOS lên các topic tương ứng. Broker đóng vai trò trung gian tiếp nhận và phân phối bản tin đến các subscriber đã đăng ký. Backend được xây dựng bằng FastAPI hoạt động như subscriber, nhận dữ liệu từ broker để lưu trữ vào cơ sở dữ liệu, xử lý nghiệp vụ và gửi thông báo khi cần thiết. Mô hình này giúp tách biệt hoàn toàn giữa thiết bị và ứng dụng phía máy chủ, cho phép mở rộng hệ thống mà không cần thay đổi chương trình trên thiết bị.

MQTT hỗ trợ ba mức chất lượng dịch vụ (Quality of Service – QoS), bao gồm QoS 0 (At most once), QoS 1 (At least once) và QoS 2 (Exactly once). Trong đề tài này, các bản tin GPS và tín hiệu SOS được truyền với mức QoS 1 nhằm đảm bảo dữ liệu được chuyển đến broker ít nhất một lần, đồng thời vẫn duy trì hiệu năng truyền phù hợp cho hệ thống theo dõi thời gian thực.

Các giao thức thay thế đã được xem xét gồm HTTP, CoAP và AMQP. HTTP sử dụng mô hình request/response nên thiết bị phải gửi yêu cầu định kỳ để cập nhật dữ liệu, làm tăng lưu lượng truyền và độ trễ. CoAP có ưu điểm về kích thước gói tin nhỏ và hoạt động trên nền UDP, tuy nhiên hệ sinh thái thư viện và mức độ phổ biến trong các ứng dụng IoT sử dụng mạng di động còn hạn chế. AMQP cung cấp nhiều tính năng nâng cao nhưng có giao thức phức tạp và chi phí tài nguyên lớn hơn so với yêu cầu của hệ thống. Vì vậy, MQTT được lựa chọn nhờ đặc tính nhẹ, hiệu quả, hỗ trợ cơ chế publish/subscribe và có hệ sinh thái thư viện phong phú trên cả ESP32 (PubSubClient) và backend Python (Paho MQTT \cite{paho_mqtt}), đáp ứng tốt yêu cầu truyền dữ liệu thời gian thực của hệ thống.

\subsection{Broker MQTT – EMQX Public Broker}
\label{subsection:3.2.2}

EMQX là MQTT Broker mã nguồn mở do EMQ Technologies phát triển, hỗ trợ các phiên bản MQTT 3.1, 3.1.1 và 5.0 \cite{emqx_broker}. Ngoài khả năng tự triển khai trên máy chủ riêng, EMQX còn cung cấp dịch vụ Public Broker tại địa chỉ \texttt{broker.emqx.io}, cho phép các thiết bị và ứng dụng kết nối để thử nghiệm và phát triển mà không cần cài đặt hoặc cấu hình máy chủ MQTT.

Trong hệ thống đề xuất, cả thiết bị ESP32 và backend FastAPI đều kết nối tới EMQX Public Broker. ESP32 thực hiện publish dữ liệu vị trí GPS và tín hiệu SOS lên các topic được quy định trước, trong khi backend đăng ký (subscribe) các topic này để tiếp nhận dữ liệu, lưu vào cơ sở dữ liệu MongoDB và xử lý các nghiệp vụ như gửi thông báo Telegram hoặc cung cấp dữ liệu cho ứng dụng di động thông qua API. Các topic được tổ chức theo cấu trúc phân cấp nhằm dễ dàng quản lý và giảm nguy cơ xung đột với các ứng dụng khác khi sử dụng broker công khai.

Việc sử dụng EMQX Public Broker giúp đơn giản hóa quá trình phát triển do không cần triển khai và quản trị hạ tầng MQTT riêng. Tuy nhiên, do đây là dịch vụ công khai nên dữ liệu được truyền qua hạ tầng dùng chung và không phù hợp cho môi trường sản xuất hoặc các hệ thống yêu cầu mức độ bảo mật cao. Khi triển khai thực tế, hệ thống có thể chuyển sang sử dụng EMQX tự triển khai hoặc các dịch vụ MQTT có hỗ trợ cơ chế xác thực và mã hóa để tăng cường tính bảo mật và độ tin cậy.

Các phương án thay thế đã được xem xét gồm HiveMQ Cloud, Eclipse Mosquitto và AWS IoT Core. HiveMQ Cloud cung cấp dịch vụ MQTT được quản lý nhưng yêu cầu đăng ký tài khoản và có giới hạn tài nguyên ở gói miễn phí. Eclipse Mosquitto cho phép tự triển khai và kiểm soát hoàn toàn hệ thống nhưng đòi hỏi quản trị máy chủ. AWS IoT Core cung cấp nhiều tính năng quản lý thiết bị và bảo mật nhưng có chi phí vận hành và cấu hình phức tạp hơn. Vì vậy, EMQX Public Broker được lựa chọn do dễ sử dụng, không yêu cầu cài đặt hạ tầng trong giai đoạn phát triển, đồng thời có thể dễ dàng thay thế bằng một MQTT Broker riêng khi hệ thống được triển khai trong môi trường thực tế.

\section{Backend}
\label{section:3.3}

\subsection{Framework FastAPI}
\label{subsection:3.3.1}

FastAPI là framework phát triển ứng dụng web dành cho Python, được xây dựng trên nền tảng Starlette và Pydantic, hỗ trợ lập trình bất đồng bộ (async/await) theo chuẩn ASGI \cite{fastapi_docs}. Framework cho phép xây dựng các RESTful API với hiệu năng cao, đồng thời tự động sinh tài liệu API theo chuẩn OpenAPI (Swagger UI), giúp việc kiểm thử và phát triển trở nên thuận tiện hơn.

Trong hệ thống đề xuất, FastAPI đảm nhiệm vai trò xây dựng backend và xử lý các yêu cầu từ ứng dụng di động cũng như thiết bị IoT. Backend đồng thời thực hiện nhiều tác vụ như tiếp nhận dữ liệu từ MQTT Broker, cung cấp REST API cho ứng dụng di động, giao tiếp với cơ sở dữ liệu MongoDB và gửi thông báo đến Telegram Bot khi phát hiện các sự kiện cảnh báo. Khả năng xử lý bất đồng bộ của FastAPI giúp hệ thống xử lý hiệu quả các tác vụ I/O, nâng cao khả năng đáp ứng khi nhiều yêu cầu được thực hiện đồng thời.

Các phương án thay thế đã được xem xét gồm Django REST Framework, Flask và Node.js/Express. Django REST Framework cung cấp nhiều tính năng tích hợp sẵn nhưng có kiến trúc tương đối lớn và không tối ưu cho các ứng dụng yêu cầu xử lý bất đồng bộ. Flask có ưu điểm là nhẹ và dễ sử dụng nhưng khả năng hỗ trợ lập trình bất đồng bộ còn hạn chế và cần bổ sung thêm nhiều thư viện để xây dựng hệ thống hoàn chỉnh. Node.js/Express có khả năng xử lý bất đồng bộ tốt, tuy nhiên việc phát triển sẽ phải sử dụng hệ sinh thái JavaScript thay vì Python, gây khó khăn trong việc tích hợp với các thư viện và mã nguồn hiện có của hệ thống. Vì vậy, FastAPI được lựa chọn nhờ hiệu năng cao, hỗ trợ lập trình bất đồng bộ, dễ phát triển và tích hợp tốt với hệ sinh thái Python.

\subsection{Cơ sở dữ liệu MongoDB}
\label{subsection:3.3.2}

MongoDB là hệ quản trị cơ sở dữ liệu NoSQL hướng tài liệu, lưu trữ dữ liệu dưới định dạng BSON (Binary JSON) \cite{mongodb_docs}. Với khả năng lưu trữ dữ liệu theo cấu trúc linh hoạt, MongoDB đặc biệt phù hợp với các hệ thống IoT khi dữ liệu được truyền dưới dạng JSON từ thiết bị. Không giống cơ sở dữ liệu quan hệ, MongoDB không yêu cầu định nghĩa trước cấu trúc bảng cố định, giúp việc bổ sung hoặc thay đổi các trường dữ liệu trong quá trình phát triển trở nên thuận tiện hơn.

Trong hệ thống đề xuất, MongoDB được sử dụng để lưu trữ thông tin tài khoản người dùng, thông tin thiết bị, dữ liệu vị trí GPS và các thông tin liên quan khác. Dữ liệu vị trí được lưu liên tục theo thời gian, phục vụ việc hiển thị vị trí hiện tại, xem lại lịch sử di chuyển và hỗ trợ các chức năng quản lý trên ứng dụng.

MongoDB được triển khai thông qua dịch vụ MongoDB Atlas, cho phép lưu trữ cơ sở dữ liệu trên nền tảng đám mây và truy cập từ backend thông qua Internet. Việc sử dụng MongoDB Atlas giúp đơn giản hóa quá trình triển khai, quản lý cơ sở dữ liệu và sao lưu dữ liệu trong giai đoạn phát triển hệ thống.

Các phương án thay thế đã được xem xét gồm PostgreSQL, InfluxDB và Firebase Realtime Database. PostgreSQL có khả năng quản lý dữ liệu quan hệ mạnh mẽ nhưng yêu cầu thiết kế lược đồ dữ liệu chặt chẽ và việc thay đổi cấu trúc dữ liệu kém linh hoạt hơn. InfluxDB được tối ưu cho dữ liệu chuỗi thời gian nhưng không phù hợp khi hệ thống cần lưu trữ đồng thời nhiều loại dữ liệu khác nhau như tài khoản người dùng, thiết bị và dữ liệu vị trí. Firebase Realtime Database dễ tích hợp với ứng dụng di động nhưng khả năng truy vấn dữ liệu phức tạp còn hạn chế. Vì vậy, MongoDB được lựa chọn nhờ tính linh hoạt, khả năng mở rộng tốt và tích hợp thuận tiện với backend Python.

\subsection{Xác thực người dùng}
\label{subsection:3.3.3}

Hệ thống sử dụng cơ chế xác thực dựa trên Bearer Authentication kết hợp với thuật toán băm mật khẩu bcrypt nhằm bảo vệ thông tin đăng nhập của người dùng. Khi người dùng đăng ký tài khoản, mật khẩu được băm bằng thuật toán bcrypt trước khi lưu vào cơ sở dữ liệu, giúp tránh lưu trữ mật khẩu dưới dạng văn bản thuần (plaintext) và tăng cường tính an toàn cho dữ liệu.

Khi người dùng đăng nhập, backend xác minh mật khẩu bằng cách so sánh với giá trị đã được băm trong cơ sở dữ liệu. Sau khi xác thực thành công, backend trả về định danh của người dùng để ứng dụng di động sử dụng trong các yêu cầu tiếp theo. Thông tin này được gửi trong trường Authorization theo chuẩn Bearer, giúp backend xác định người dùng tương ứng và xử lý các yêu cầu truy cập tài nguyên.

Cơ chế xác thực này có ưu điểm là đơn giản, không cần duy trì trạng thái phiên làm việc (stateless) và phù hợp với quy mô của hệ thống trong giai đoạn phát triển và thử nghiệm. Trong tương lai, khi hệ thống được triển khai thực tế với số lượng người dùng lớn hơn và yêu cầu bảo mật cao hơn, cơ chế xác thực có thể được mở rộng bằng các giải pháp như JWT hoặc OAuth 2.0 nhằm tăng cường tính bảo mật và khả năng quản lý phiên đăng nhập.

\section{Ứng dụng di động}
\label{section:3.4}

\subsection{Framework Flutter}
\label{subsection:3.4.1}

Flutter là framework phát triển ứng dụng đa nền tảng do Google phát triển, sử dụng ngôn ngữ lập trình Dart và biên dịch trực tiếp thành mã máy (Ahead-of-Time Compilation – AOT) khi phát hành ứng dụng \cite{flutter_docs}. Flutter sử dụng engine đồ họa riêng để kết xuất giao diện, cho phép xây dựng giao diện thống nhất trên nhiều nền tảng mà không phụ thuộc vào các thành phần giao diện của hệ điều hành.

Trong hệ thống đề xuất, Flutter được sử dụng để phát triển ứng dụng dành cho phụ huynh trên cả Android và iOS từ cùng một mã nguồn. Ứng dụng thực hiện các chức năng như đăng nhập, quản lý thiết bị, theo dõi vị trí của trẻ trên bản đồ, xem lịch sử di chuyển và nhận các thông báo cảnh báo từ hệ thống. Việc sử dụng một codebase duy nhất giúp giảm thời gian phát triển, đồng thời thuận tiện cho việc bảo trì và mở rộng ứng dụng trong tương lai.

Ứng dụng sử dụng thư viện \textbf{flutter\_map} \cite{flutter_map_pkg} để hiển thị bản đồ tương tác và thư viện \textbf{web\_socket\_channel} để duy trì kết nối WebSocket với backend nhằm nhận dữ liệu vị trí theo thời gian thực. Việc quản lý trạng thái của ứng dụng được thực hiện thông qua Provider, giúp tổ chức mã nguồn rõ ràng và phù hợp với quy mô của hệ thống.

Các phương án thay thế đã được xem xét gồm React Native, phát triển ứng dụng native riêng cho Android và iOS, và Kotlin Multiplatform Mobile. React Native cho phép phát triển đa nền tảng nhưng phụ thuộc vào JavaScript và hệ thống bridge để giao tiếp với các thành phần native. Phát triển native bằng Kotlin và Swift mang lại khả năng khai thác tối đa từng nền tảng nhưng làm tăng khối lượng phát triển và bảo trì do phải xây dựng hai ứng dụng riêng biệt. Kotlin Multiplatform Mobile là một giải pháp đa nền tảng mới nhưng hệ sinh thái và số lượng thư viện hỗ trợ còn hạn chế. Vì vậy, Flutter được lựa chọn nhờ khả năng phát triển đa nền tảng từ một mã nguồn duy nhất, hiệu năng tốt, giao diện thống nhất và hệ sinh thái thư viện phong phú.

\subsection{Bản đồ OpenStreetMap}
\label{subsection:3.4.2}

OpenStreetMap (OSM) là dự án bản đồ mã nguồn mở được xây dựng và cập nhật bởi cộng đồng trên toàn thế giới, cung cấp dữ liệu bản đồ miễn phí phục vụ nhiều lĩnh vực khác nhau \cite{osm_wiki}. Thông qua thư viện flutter\_map, ứng dụng có thể hiển thị bản đồ, đánh dấu vị trí thiết bị GPS và hỗ trợ người dùng theo dõi vị trí trực quan ngay trên thiết bị di động.

Trong hệ thống đề xuất, dữ liệu tọa độ GPS do thiết bị gửi về được hiển thị trên nền bản đồ OpenStreetMap. Ứng dụng sử dụng các marker để biểu diễn vị trí hiện tại của trẻ, đồng thời hỗ trợ cập nhật vị trí theo thời gian thực thông qua dữ liệu nhận từ backend. Việc sử dụng OpenStreetMap kết hợp với flutter\_map giúp hệ thống dễ dàng tích hợp chức năng bản đồ mà không cần phụ thuộc vào các SDK độc quyền.

Phương án thay thế phổ biến là Google Maps SDK. Giải pháp này cung cấp nhiều tính năng nâng cao như chỉ đường, tìm kiếm địa điểm và dữ liệu bản đồ phong phú, tuy nhiên yêu cầu cấu hình API Key, sử dụng tài khoản Google Cloud và có thể phát sinh chi phí khi vượt quá giới hạn miễn phí. Đối với hệ thống trong phạm vi nghiên cứu và phát triển nguyên mẫu, các chức năng chính chỉ bao gồm hiển thị vị trí và theo dõi thiết bị trên bản đồ. Vì vậy, OpenStreetMap được lựa chọn nhờ là nền tảng mã nguồn mở, miễn phí sử dụng, dễ tích hợp với Flutter và đáp ứng đầy đủ yêu cầu của hệ thống.

\section{Thuật toán Haversine}
\label{section:3.5}

Thuật toán Haversine được sử dụng để tính khoảng cách đường chim bay giữa hai điểm trên bề mặt Trái Đất dựa trên tọa độ vĩ độ và kinh độ của chúng \cite{haversine_formula}. Thuật toán giả định Trái Đất có dạng hình cầu và được sử dụng rộng rãi trong các ứng dụng định vị, bản đồ và hệ thống theo dõi vị trí.

Công thức tính khoảng cách được biểu diễn như sau:

\begin{equation}
a = \sin^2\!\left(\frac{\Delta\varphi}{2}\right) + \cos\varphi_1 \cdot \cos\varphi_2 \cdot \sin^2\!\left(\frac{\Delta\lambda}{2}\right)
\label{eq:haversine_a}
\end{equation}

\begin{equation}
d = 2R \cdot \arctan2\!\left(\sqrt{a},\, \sqrt{1-a}\right)
\label{eq:haversine_d}
\end{equation}

\noindent trong đó $\varphi_1$ và $\varphi_2$ lần lượt là vĩ độ của hai điểm (đơn vị radian); $\Delta\varphi$ và $\Delta\lambda$ là chênh lệch vĩ độ và kinh độ giữa hai điểm; $R$ là bán kính trung bình của Trái Đất; và $d$ là khoảng cách giữa hai vị trí cần tính.

Trong hệ thống đề xuất, thuật toán Haversine được sử dụng tại backend để kiểm tra thiết bị có nằm trong vùng an toàn (geofence) hay không. Mỗi khi nhận được tọa độ GPS mới từ thiết bị, backend sẽ tính khoảng cách từ vị trí hiện tại đến tâm của từng vùng geofence đã được thiết lập. Nếu khoảng cách vượt quá bán kính của vùng an toàn, hệ thống sẽ xác định thiết bị đã ra khỏi geofence và thực hiện các chức năng cảnh báo như gửi thông báo đến ứng dụng di động và Telegram Bot.

Các phương án thay thế đã được xem xét gồm Google Geofencing API, PostGIS và công thức Vincenty. Google Geofencing API cung cấp sẵn chức năng kiểm tra geofence nhưng phụ thuộc vào dịch vụ của bên thứ ba và có thể phát sinh chi phí sử dụng. PostGIS hỗ trợ mạnh mẽ các truy vấn không gian nhưng yêu cầu sử dụng PostgreSQL thay cho MongoDB. Công thức Vincenty có độ chính xác cao hơn do sử dụng mô hình ellipsoid của Trái Đất, tuy nhiên việc tính toán phức tạp hơn và không cần thiết đối với phạm vi theo dõi của hệ thống. Vì vậy, thuật toán Haversine được lựa chọn nhờ cách triển khai đơn giản, không phụ thuộc dịch vụ bên ngoài và có độ chính xác phù hợp với yêu cầu xác định vùng an toàn của hệ thống.

Kết chương: Chương 3 đã giới thiệu và phân tích lý do lựa chọn toàn bộ công nghệ trong hệ thống. Sự kết hợp giữa ESP32 và SIM7600 giải quyết yêu cầu về phần cứng nhúng tự chủ với kết nối 4G LTE ổn định; MQTT qua EMQX public broker đảm bảo truyền dữ liệu tin cậy trong điều kiện mạng di động; FastAPI và MongoDB cung cấp nền tảng backend async linh hoạt; Flutter và OSM cho phép xây dựng ứng dụng di động đa nền tảng không phụ thuộc dịch vụ có phí; và thuật toán Haversine đáp ứng nhu cầu geofencing hoàn toàn tự chủ. Trên nền tảng công nghệ này, Chương 4 sẽ trình bày chi tiết quá trình thiết kế và triển khai từng thành phần của hệ thống.

\end{document}
